\documentclass[../../main2.tex]{subfiles}

\begin{document}

\chapter{通道一：观念的病毒（无压扩增）}
\label{ch:meme}

\begin{motivation}
	上一章遗留的问题：生成树的第一枝——不加任何约束的传播长什么样？什么信息会\textbf{自己}传播？\\
	\textbf{核心动机}：定义模因（meme），确立它的关键指标是传播强度而非正确性，并把宗教改革、谣言、短视频这三个相隔五百年的现场，放进同一套账。\\
	\textbf{扩增模式}：无压——复制不消耗原体，上限在接收方。
\end{motivation}

\begin{logicmap}
\begin{tikzpicture}[node distance=0.85cm, every node/.style={draw, rectangle, rounded corners, align=center, font=\small}]
\node (q) {什么信息会自己传播？};
\node (d) [below=of q] {模因：不耗原体的复制};
\node (s) [below=of d] {强度 > 正确性};
\node (l) [below=of s] {语言：压缩层};
\node (c) [below=of l] {历史现场：宗教改革与谣言};
\node (n) [below=of c] {传播要花钱呢？（抛给第 12 章）};
\draw[-{Stealth}] (q) -- (d);
\draw[-{Stealth}] (d) -- (s);
\draw[-{Stealth}] (s) -- (l);
\draw[-{Stealth}] (l) -- (c);
\draw[-{Stealth}] (c) -- (n);
\end{tikzpicture}
\end{logicmap}

\section{模因：不耗原体的复制}

\begin{readingnote}
	你有没有想过，为什么谣言永远比辟谣跑得快？\\
	不是造谣的人多——是\textbf{两者不在同一条起跑线上}。
\end{readingnote}

先给定义。模因（meme）：一段\textbf{能自我复制的信息}——复制不消耗原体，上限只在接收方的注意力和理解力。笑话、口号、谣言、表情包、习俗，都是模因。三个可操作的判据：

\begin{itemize}
	\item \textbf{可复制}：换个载体内容不丢（故事从口到书到视频，还是那个故事）；
	\item \textbf{不耗原体}：你讲完笑话，你还有笑话；
	\item \textbf{以接收方为瓶颈}：没人听，就断了——扩增的天花板是「多少人接得住」。
\end{itemize}

\begin{tcolorbox}[colback=green!5, title=\textit{生活类比}]
	模因像兔子不像金币：金币分一块少一块，兔子一窝生一窝。所以「观念的病毒」这个外号很准——病毒正是靠复制不耗原体，才能从一个人滚到几十亿人。
\end{tcolorbox}

跨人群的模因，你手里就攥着几个：

\begin{itemize}
	\item \textbf{学生}：「内卷」「上岸」——两个词概括一整套处境，所以传得动。
	\item \textbf{上班族}：「赋能」「抓手」——黑话是职场模因，先传起来，再想意思。
	\item \textbf{刷手机的你}：那个转疯了的表情包——没人考证过它的出处，没人需要考证。
	\item \textbf{古代}：「苍天已死，黄天当立」——八个字动员了百万流民。模因从来是起义的第一件武器。
\end{itemize}

\paragraph*{逻辑衔接}
	模因认出来了：不耗原体的复制。但同样一段信息，有的胎死腹中，有的传遍全网——谁在决定生死？下一节：传播强度。

\section{强度压倒正确性}

\begin{readingnote}
	你有没有想过，「震惊体」标题的写作水平，为什么常常好过正规新闻？\\
	因为它的优化目标不是「真」，是\textbf{转}。
\end{readingnote}

模因的传播不问对错，只问强度（intensity）。三个燃料，按点火顺序排：

\begin{enumerate}
	\item \textbf{情绪激活}：愤怒、恐惧、优越感——情绪是转发的燃料。让你「气得想让别人也看看」的信息，赢在起跑线。
	\item \textbf{认知成本}：一句话能懂的打败三页纸的。「多喝热水」传播成本趋近于零。
	\item \textbf{身份认同}：「我们 vs 他们」是模因的放大器——转它等于亮明身份。
\end{enumerate}

为什么强度压倒正确性？因为\textbf{验证是接收方的成本，情绪是传播方的燃料}。辟谣要读完、要查证、要承认自己错了——三道成本全是负激励；造谣只需要「哇」。

\begin{questionbox}
	\textit{现在你可能想反驳：「照这么说，真理永远输给谣言？那人类怎么进步的？」}\\
	\textbf{长期看，真理有制度后援。}科学期刊、判决书、教科书——都是给「验证」发工资的机构。模因通道内谣言赢，跨通道看，经济和法律会来纠偏（这正是后面几章的戏）。乐观有理由，但别指望单靠转发真理就能赢。
\end{questionbox}

\begin{warningbox}
	\textbf{平台动机插一脚}（第 10 章那个问题的续集）：算法为什么专推你爱看的？因为平台的收益是停留时长，而\textbf{高强度模因最抓停留}。于是「情绪优先」的模因筛选，被商业动机放大了几十倍。谣言跑得快，一半是人性，一半是 KPI。
\end{warningbox}

\paragraph*{逻辑衔接}
	强度定生死。但高强度的模因传久了，会把自己压进语言里——下一节：语言是模因的最高频压缩层。

\section{语言：压缩层}

\begin{readingnote}
	你有没有想过，为什么「成语」四个字能顶一个故事？\\
	压缩到极致的模因，传得最远。
\end{readingnote}

语言是模因的高速公路加收费站：成语「守株待兔」四个字封装了一个寓言，学生学一遍终身会用；梗「真香」两个字封装了一整套打脸叙事。\textbf{压缩率越高，传播成本越低}——这就是黑话、缩写、梗图永远在赢的原因。

反过来的代价也要认：压缩必有丢失。「内卷」传遍全网之后，原义（人类学里的精细耕作）几乎没人记得。模因越成功，越容易变成空壳——这是传播强度和语义保真之间的交易。

\paragraph*{逻辑衔接}
	机制齐了：模因靠强度扩散，靠压缩落地。现在拉三个真实现场，把这三件事对号入座——其中两个，隔了五百年。

\section{三个现场：宗教改革、谣言与茧房}

\textbf{现场一：宗教改革（1517）}。马丁·路德的《九十五条论纲》是模因教科书。它赢在三处全中：情绪激活（赎罪券的愤怒）、认知成本（论纲用德语写、可辩论、可传抄）、身份认同（德意志 vs 罗马）。印刷术是这次的「算法」——复制成本暴跌，模因无压扩增的上限被抬高了一个数量级。\lowfreq{（把印刷术类比算法是半严谨类比：印刷没有个性化推荐，但都改变了复制成本。）}三十年战争死了几百万人——一个模因的下游账单。

\textbf{现场二：流行病谣言}。「板蓝根防病毒」「熏醋消毒」——每次疫情都重演一遍。情绪（怕死）拉满，认知成本（一句话）趋零，身份（「我是懂防护的人」）拉满。辟谣要写论文。三比零，谣言赢。

\textbf{现场三：信息茧房}。茧房的白话版第 7 章讲过：算法按你的点击喂料。放进本章的账里：茧房是\textbf{模因通道被个性化分发改造后}的产物——同一个模因池，每个人被喂了不同的切片，于是「同一场舆论」裂成一百个平行现场。宏观的「舆论」，在通道层面就散架了。

\begin{questionbox}
	\textit{现在你可能想反驳：「王安石变法失败，你总不能也赖模因吧？」}\\
	\textbf{还真能算上一份。}新法的每一条都先要当模因传——「青苗法」在士大夫嘴里被传成「官府放高利贷」，一旦这个压缩版定型，朝廷再解释就没用了。观念战输了，执行就变形。动机好、渠道被劫持——第 13 章政治里还有他一笔账。
\end{questionbox}

\paragraph*{逻辑衔接}
	模因通道的账能估了：找一个具体观念，问它「拿掉会怎样」（无你检验），再看强度和压缩。但模因自己免费，它的\textbf{传播}常常不免费——印书要钱，买量要钱，办报要钱。一旦传播必须付费，还只是模因吗？第 12 章：加入稀缺约束，经济被逼出来。

\begin{thinkbox}
\begin{enumerate}
	\item 挑一个这周刷到、转给过朋友的段子或观点：拆它的三个燃料（情绪、认知成本、身份）——它是靠哪个赢的？把标题改成「冷静客观版」，它还会传吗？
	\item \textbf{反事实}：如果 1517 年没有印刷术，《九十五条论纲》还会引发宗教改革吗？换成口耳相传，这场观念的病毒要传多少年？改写欧洲史哪几页？
	\item \textbf{反事实}：如果辟谣和造谣的传播成本一样（比如平台强制「转发前先看三秒事实核查」），谣言的强度优势还剩多少？你的信息茧房会变薄吗？
\end{enumerate}
\end{thinkbox}

\end{document}
